% ---------------------------------------------------------------------
%  CHƯƠNG 5 — KẾT QUẢ, ĐÁNH GIÁ VÀ KẾT LUẬN
% ---------------------------------------------------------------------
\chapter{Kết quả, đánh giá và kết luận}
\label{ch:ketqua}

\begin{muctieu}
Chương này đối chiếu sản phẩm thu được với từng mục tiêu đặt ra ở
Chương~\ref{ch:modau}, trình bày kết quả kiểm thử tự động, cách đo và kết
quả đo hiệu năng, bàn luận về những điểm hệ
thống làm được và chưa làm được, rồi tổng kết đồ án, nêu đóng góp chính và
các hướng phát triển tiếp theo.
\end{muctieu}

\section{Đối chiếu với mục tiêu ban đầu}

\begin{table}[H]
\centering
\caption{Mức độ hoàn thành từng mục tiêu nêu ở Mục~\ref{sec:muctieu}.}
\label{tab:doichieumuctieu}
\small
\begin{tabularx}{\textwidth}{@{}>{\bfseries}c >{\raggedright\arraybackslash}p{2.2cm} X@{}}
\toprule
Mục tiêu & Kết quả & Bằng chứng \\
\midrule
MT1 & Đạt & Quy trình đăng ký và nhận diện chạy trọn vẹn trên ONNX Runtime
  bản CPU, không cần GPU. Thời gian xử lý đo được nêu ở
  Bảng~\ref{tab:hieunang} \\
MT2 & Đạt & Sáu lớp phòng vệ ở Bảng~\ref{tab:phongve} đều đã cài đặt và có
  kiểm thử tự động tương ứng \\
MT3 & Đạt & \FsoEndpoint{} endpoint trên \FsoBang{} bảng dữ liệu phủ hết
  vòng đời từ hồ sơ tới báo cáo, có nhật ký sửa điểm danh, nhật ký hệ thống,
  thông báo cho phụ huynh và xuất Excel \\
MT4 & Đạt & Ma trận vai trò và chức năng ở Bảng~\ref{tab:phanquyen}, gồm cả
  việc tách văn phòng Khoa với lãnh đạo Khoa, được cài đặt tập trung và kiểm
  chứng bằng nhóm kiểm thử phân quyền \\
MT5 & Đạt & Một mã nguồn Flutter với \FsoManHinh{} màn hình chạy trên Web,
  Android, iOS; \FsoTest{} ca kiểm thử tự động chạy trên PostgreSQL thật \\
\bottomrule
\end{tabularx}
\end{table}

{\small
\begin{xltabular}{\textwidth}{@{}>{\bfseries\raggedright\arraybackslash\leavevmode}p{2.0cm} >{\raggedright\arraybackslash\leavevmode}p{5.6cm} >{\raggedright\arraybackslash\leavevmode}X@{}}
\caption{Giải thích chi tiết Bảng~\ref{tab:doichieumuctieu}: mỗi mục tiêu được kiểm chứng ở đâu.}
\label{tab:doichieumuctieu-gt}\\
\toprule
Mục tiêu & Yêu cầu tóm tắt & Kiểm chứng tại \\
\midrule
\endfirsthead
\toprule
Mục tiêu & Yêu cầu tóm tắt & Kiểm chứng tại \\
\midrule
\endhead
\bottomrule
\endfoot
MT1 & Nhận diện trọn vẹn trên máy chỉ có CPU, đủ nhanh để dùng trong lớp & Bảng~\ref{tab:hieunang} và Mục~\ref{sec:hieunang}; khung đông người còn chậm, ghi ở phần hạn chế \\
MT2 & Chống ảnh in hoặc màn hình, chống nhận nhầm do một khung xấu, chống điểm danh từ xa & Bảng~\ref{tab:phongve} và các tệp kiểm thử tương ứng ở Bảng~\ref{tab:kiemthu}; hiệu quả chặn ngoài thực tế chưa đo \\
MT3 & Số hoá trọn nghiệp vụ chuyên cần & Bảng~\ref{tab:api} và Bảng~\ref{tab:luocdo} \\
MT4 & Phân quyền đúng mô hình tổ chức & Bảng~\ref{tab:phanquyen} và các tệp kiểm thử phân quyền \\
MT5 & Một ứng dụng cho ba nền tảng, có kiểm thử tự động & Bảng~\ref{tab:manhinh} và Bảng~\ref{tab:kiemthu} \\
Cột Kết quả & Đạt nghĩa là có bằng chứng kiểm chứng được nêu ở cột cuối & Những điểm còn thiếu ghi ở mục hạn chế của chương này \\
\end{xltabular}}

\section{Kết quả kiểm thử tự động}

Bộ kiểm thử viết bằng pytest gồm \FsoTest{} ca trong \FsoFileTest{} tệp, khoảng
\FsoDongTest{} dòng mã kiểm thử. Bảng~\ref{tab:kiemthu} liệt kê từng nhóm.

{\scriptsize
\begin{xltabular}{\textwidth}{@{}>{\raggedright\arraybackslash\leavevmode}p{5.6cm} >{\centering\arraybackslash}p{1.3cm} X@{}}
\caption{Các nhóm kiểm thử tự động của phía máy chủ.}
\label{tab:kiemthu}\\
\toprule
Tệp kiểm thử & Số ca & Nội dung kiểm \\
\midrule
\endfirsthead
\toprule
Tệp kiểm thử & Số ca & Nội dung kiểm \\
\midrule
\endhead
\bottomrule
\endfoot
\BangTest
\end{xltabular}}

{\footnotesize
\begin{xltabular}{\textwidth}{@{}>{\bfseries\raggedright\arraybackslash\leavevmode}p{2.6cm} >{\raggedright\arraybackslash\leavevmode}p{6.4cm} >{\raggedright\arraybackslash\leavevmode}X@{}}
\caption{Giải thích chi tiết Bảng~\ref{tab:kiemthu}: các tệp kiểm thử theo nhóm rủi ro.}
\label{tab:kiemthu-gt}\\
\toprule
Nhóm & Tệp & Rủi ro được chặn \\
\midrule
\endfirsthead
\toprule
Nhóm & Tệp & Rủi ro được chặn \\
\midrule
\endhead
\bottomrule
\endfoot
Xác thực, phiên đăng nhập & \code{test\_auth}, \code{test\_tokens\_and\_media}, \code{test\_password\_reset}, \code{test\_change\_password}, \code{test\_manager\_reset\_password} & Lộ tài khoản, token bị đánh cắp được dùng lại \\
Phân quyền & \code{test\_authorization}, \code{test\_unit\_roles}, \code{test\_admin\_management}, \code{test\_admin\_crud} & Xem hoặc sửa dữ liệu ngoài phạm vi của vai trò \\
Nhận diện, chống gian lận & \code{test\_antispoof}, \code{test\_confirm\_frames}, \code{test\_location\_checkin}, \code{test\_checkin\_qr}, \code{test\_consent}, \code{test\_face\_approval} & Điểm danh hộ, xử lý dữ liệu sinh trắc khi chưa có đồng ý \\
Vòng đời điểm danh & \code{test\_open\_session\_rules}, \code{test\_session\_lifecycle}, \code{test\_session\_detail\_and\_records}, \code{test\_rooms}, \code{test\_week\_schedule} & Mở phiên sai giờ, xếp trùng phòng, thiếu bản ghi vắng \\
Báo cáo, thông báo, nhật ký & \code{test\_reports}, \code{test\_reports\_export\_and\_warnings}, \code{test\_my\_attendance\_detail}, \code{test\_parent\_notifications}, \code{test\_audit\_log}, \code{test\_evidence\_cleanup} & Số liệu sai, quên báo phụ huynh, thiếu dấu vết thao tác \\
Lỗi đã gặp & \code{test\_flow\_fixes} & Lỗi đã sửa quay trở lại \\
Cách đếm & -- & Số ca là số hàm kiểm thử trong tệp, đếm tự động; kiểm thử của ứng dụng Flutter chạy riêng bằng \code{flutter test}, không tính vào bảng \\
\end{xltabular}}

Ba quyết định về cách kiểm thử đáng được nêu rõ:

\begin{itemize}[leftmargin=1.6em, itemsep=3pt]
  \item \textbf{Chạy trên PostgreSQL thật, không dùng cơ sở dữ liệu nhúng.}
        Kiểu \code{Vector(\FchieuVector)} và toán tử khoảng cách của pgvector
        không có tương đương trong SQLite, nên nếu kiểm thử bằng cơ sở dữ
        liệu nhúng thì đúng phần quan trọng nhất lại không được kiểm.
  \item \textbf{Mỗi ca chạy trong một giao dịch rồi quay lui.} Nhờ vậy các ca
        không ảnh hưởng lẫn nhau mà không phải dựng lại lược đồ mỗi lần.
  \item \textbf{Giảm số vòng băm mật khẩu khi kiểm thử.} Bcrypt cố tình chậm \cite{Provos1999Bcrypt};
        hạ số vòng trong môi trường kiểm thử giúp bộ kiểm thử chạy nhanh mà
        không đổi logic nghiệp vụ.
\end{itemize}

Kết quả lần chạy gần nhất: toàn bộ \FsoTest{} ca đều đạt, thời gian chạy
39,56 giây (\code{278 passed in 39.56s}), thực hiện ngày 12/11/2026 trên
PostgreSQL 16 chạy trong Docker.

\section{Đánh giá hiệu năng}
\label{sec:hieunang}

\subsection{Cách đo}

Hiệu năng được đo ở ba điểm tách biệt: thời gian phát hiện và trích đặc trưng
một khung hình, thời gian truy vấn so khớp trong cơ sở dữ liệu, và thời gian
xử lý trọn vẹn một yêu cầu nhận diện tính từ lúc máy chủ nhận khung hình. Mỗi
phép đo lặp lại \HnSoLan{} lần, bỏ lần chạy đầu tiên vì lần đó còn phải nạp mô
hình, rồi lấy trung vị thay cho trung bình để giảm ảnh hưởng của các lần chạy
bất thường.

Phép đo do công cụ \code{backend/tools/do\_hieu\_nang.py} thực hiện ngày
\HnNgayDo{} trên máy thử nghiệm ở Bảng~\ref{tab:moitruong}. Vì chưa có ảnh
chụp bằng thiết bị sẽ dùng khi bảo vệ, phép đo dùng ảnh mẫu đi kèm thư viện
InsightFace: ảnh nhóm \code{t1.jpg} kích thước 1280$\times$886 có \HnSoMat{}
khuôn mặt, và một khung 640$\times$480 cắt từ chính ảnh đó, chứa đúng một khuôn
mặt. Phép truy vấn chạy trên cơ sở dữ liệu thử nghiệm có \HnSoVector{} vector
khuôn mặt. Công cụ ghi thẳng kết quả vào bảng dưới đây và vào các con số trong
phần nhận xét, nên số trong báo cáo là số máy in ra chứ không phải số chép tay.

\begin{table}[H]
\centering
\caption{Kết quả đo hiệu năng trên máy thử nghiệm.}
\label{tab:hieunang}
\small
\begin{tabularx}{\textwidth}{@{}X >{\centering\arraybackslash}p{3.0cm} >{\centering\arraybackslash}p{3.0cm}@{}}
\toprule
Phép đo & Trung vị & Lớn nhất \\
\midrule
\BangHieuNang
\bottomrule
\end{tabularx}
\end{table}

{\small
\begin{xltabular}{\textwidth}{@{}>{\bfseries\raggedright\arraybackslash\leavevmode}p{2.8cm} >{\raggedright\arraybackslash\leavevmode}X >{\raggedright\arraybackslash\leavevmode}p{4.6cm}@{}}
\caption{Giải thích chi tiết Bảng~\ref{tab:hieunang}: mỗi phép đo gồm những gì.}
\label{tab:hieunang-gt}\\
\toprule
Phép đo & Gồm những bước & Đầu vào \\
\midrule
\endfirsthead
\toprule
Phép đo & Gồm những bước & Đầu vào \\
\midrule
\endhead
\bottomrule
\endfoot
Một khuôn mặt & Chạy cả bộ mô hình đang nạp: phát hiện, ba mô hình phụ, trích đặc trưng & Khung 640$\times$480 cắt từ ảnh mẫu, đúng một khuôn mặt \\
Nhiều khuôn mặt & Như trên, trên cả ảnh nhóm & Ảnh mẫu \code{t1.jpg} 1280$\times$886, \HnSoMat{} khuôn mặt \\
Truy vấn so khớp & Một câu SQL tìm vector gần nhất bằng toán tử cosine của pgvector & Cơ sở dữ liệu thử nghiệm có \HnSoVector{} vector; vector truy vấn ngẫu nhiên \\
Chống giả mạo & Cắt vùng mặt mở rộng 2,7 lần, đưa về 80$\times$80, chạy MiniFASNetV2 & Khuôn mặt trong khung một khuôn mặt \\
Trọn vẹn một yêu cầu & Hàm \code{recognize\_frame} mà endpoint quét camera gọi: giải mã ảnh, chạy mô hình, chống giả mạo, so khớp trong phạm vi lớp; không gồm thời gian truyền mạng & Khung một khuôn mặt, lớp học phần đầu tiên trong cơ sở dữ liệu \\
Trung vị, lớn nhất & Trung vị và lần chậm nhất trong \HnSoLan{} lần đo, đã bỏ lần đầu vì lần đó còn nạp mô hình & -- \\
\end{xltabular}}


\subsection{Nhận xét về hiệu năng}

Gần như toàn bộ thời gian nằm ở phần chạy mô hình. Phát hiện và trích đặc
trưng một khuôn mặt mất \HnMotMat{}\,ms, trong khi kiểm tra chống giả mạo chỉ
mất \HnChongGiaMao{}\,ms và truy vấn so khớp \HnTruyVan{}\,ms. Con số truy vấn
nhỏ một phần vì cơ sở dữ liệu thử nghiệm ít dữ liệu; khi chạy thật, mỗi lần
so khớp chỉ xét vector của sinh viên trong lớp đang điểm danh, nên số vector
phải so vẫn chỉ cỡ sĩ số một lớp dù toàn trường có bao nhiêu sinh viên.

Mỗi khuôn mặt thêm vào khung hình tốn thêm khoảng \HnMoiMatThem{}\,ms: khung
\HnSoMat{} khuôn mặt mất \HnNhieuMat{}\,ms. Màn quét camera của giảng viên chụp
2 giây một lần và chỉ gửi khung mới khi khung trước đã xử lý xong, nên khung
đông người làm giãn nhịp quét chứ không làm dồn yêu cầu. Theo cách tính này,
một khung chụp trọn một lớp bốn mươi người sẽ mất nhiều giây; đây là ước
lượng, chưa đo trên lớp thật.

Phép đo cũng chỉ ra một chỗ tối ưu được ngay. Hệ thống đang nạp cả năm mô hình
của bộ \code{buffalo\_l}, trong đó ba mô hình phụ không được dùng tới. Đo lại
trên cùng ảnh với cấu hình chỉ nạp phát hiện và trích đặc trưng: khung một
khuôn mặt còn \HnGonMotMat{}\,ms, giảm \HnGiamMotMat\%; khung \HnSoMat{} khuôn
mặt còn \HnGonNhieuMat{}\,ms, giảm \HnGiamNhieuMat\%. Vector đặc trưng thu được
giống hệt cấu hình đầy đủ. Thay đổi này chưa được áp dụng vào hệ thống.

\section{Bàn luận và bài học rút ra}

\textbf{Phần mô hình không phải phần khó nhất.} Việc ghép SCRFD và ArcFace
vào hệ thống mất ít công hơn nhiều so với việc thiết kế đúng các quy tắc bao
quanh nó: ai được xem ảnh bằng chứng, phiên đóng lúc nào, sửa điểm danh thì
lưu vết ra sao, sinh viên rút lại dữ liệu sinh trắc thì xoá những gì. Đây là
kết luận đáng chú ý nhất rút ra từ đồ án.

\textbf{Thu hẹp không gian tìm kiếm là quyết định hiệu quả nhất.} Chỉ so khớp
trong phạm vi lớp học phần vừa giảm chi phí tính toán, vừa cắt giảm mạnh xác
suất nhận nhầm, vì hai người ngẫu nhiên giống nhau trong một lớp bốn mươi
người hiếm hơn nhiều so với trong cơ sở dữ liệu toàn trường.

\textbf{Sai số GPS trong nhà là rào cản thực tế.} Trong quá trình thử, việc
đối chiếu vị trí phải nới nhiều lần mới dùng được: trừ sai số thiết bị báo lên
trước khi so với bán kính, chỉ từ chối khi sai số vượt \FsaiSoGps{}\,m, và
kể cả khi phòng học có toạ độ riêng thì bán kính cũng không nhỏ hơn
\FbanKinhPhong\,m. Vì GPS không bao giờ đủ chính xác để phân biệt hai phòng
cạnh nhau, những buổi cần chắc chắn được giao cho mã QR xoay vòng. Đây là ví dụ cho thấy một tính
năng đúng về lý thuyết vẫn có thể không dùng được nếu bỏ qua đặc tính của
thiết bị thật.

\textbf{Phân quyền phải theo cách tổ chức làm việc thật.} Thiết kế ban đầu
gộp mọi cán bộ Khoa vào một vai trò vừa nhập liệu vừa duyệt nghỉ có phép. Khi
đối chiếu với cách Khoa vận hành, vai trò này được tách thành văn phòng Khoa
và lãnh đạo Khoa. Nhờ phạm vi dữ liệu đã được cài đặt tập trung, việc tách
chỉ cần thêm một vai trò vào nhóm vai trò cấp Khoa và chỉnh danh sách vai trò
ở các endpoint sửa dữ liệu, không phải viết lại bộ lọc nào.

\textbf{Kiểm thử tự động không thay được việc dùng thử thật.} Phiên bản đầu
của mã QR vượt qua toàn bộ kiểm thử nhưng chỉ dùng được từ bên trong ứng dụng:
ảnh QR chứa một chuỗi mã, phải vào đúng mục quét mã mới đọc được, còn camera
thường của điện thoại không biết làm gì với chuỗi đó. Thiếu sót này không lộ
ra ở kiểm thử mà ở lúc chạy thử theo thao tác của người dùng, và là lý do mã
QR được đổi sang dạng đường link. Lớp chống giả mạo cũng gặp tình huống tương
tự: kiểm thử dùng mô hình giả nên đều đạt, trong khi mô hình thật lại chấm ảnh
người thật là giả mạo vì mã nguồn chia điểm ảnh cho 255, khác với cách mô hình
được huấn luyện \cite{Minivision2020Silent}. Lỗi chỉ lộ ra khi chạy với tệp mô
hình và ảnh thật.

\textbf{Ràng buộc nên đặt ở tầng thấp nhất có thể.} Những quy tắc như mỗi lịch
học chỉ có một phiên trong ngày được đặt thành ràng buộc của cơ sở dữ liệu, nên
hai yêu cầu gửi đồng thời cũng không phá được. Nếu chỉ kiểm tra bằng lệnh điều
kiện ở tầng ứng dụng thì lỗi kiểu này rất khó tái hiện và rất dễ lọt.

\textbf{Viết kiểm thử ngay khi cài đặt là lựa chọn tiết kiệm.} Nhóm kiểm thử về
các lỗi luồng nghiệp vụ đã phát hiện là bằng chứng rõ nhất: mỗi lỗi tìm ra đều
được cố định lại bằng một ca kiểm thử, nhờ đó những lần sửa sau không làm hỏng
lại thứ đã chạy đúng.

\section{Hạn chế của hệ thống}

\begin{enumerate}[leftmargin=1.8em, itemsep=3pt]
  \item \textbf{Chưa đo độ chính xác và khả năng chống giả mạo trên dữ liệu
        thật.} Hai phép đo này cần ảnh của nhiều người thật chụp ở nhiều điều
        kiện, có sự đồng ý của người được chụp, và cần thử trực tiếp bằng ảnh
        in, màn hình và video trước camera. Đồ án mới kiểm chứng phần này bằng
        kiểm thử tự động và chạy mô hình trên ảnh mẫu, nên chưa có con số về tỉ
        lệ nhận đúng hay tỉ lệ chặn giả mạo.
  \item \textbf{Khung hình đông người chậm trên CPU.} Mỗi khuôn mặt thêm vào
        tốn khoảng \HnMoiMatThem{}\,ms (Mục~\ref{sec:hieunang}), nên quét một
        khung có cả lớp sẽ mất nhiều giây.
  \item \textbf{Ngưỡng nhận diện cố định.} Hệ thống dùng một ngưỡng chung
        $\tau = \FnguongNhanDien$ cho mọi người, trong khi ngưỡng tối ưu có
        thể khác nhau tuỳ chất lượng ảnh đăng ký của từng sinh viên.
  \item \textbf{Chống giả mạo chỉ dựa trên ảnh tĩnh.} Mô hình bị động không
        chống được những kiểu tấn công tinh vi hơn như mặt nạ ba chiều \cite{Ramachandra2017Survey}.
  \item \textbf{Bộ đếm xác nhận đa khung hình không bền vững.} Bộ đếm nằm
        trong bộ nhớ tiến trình, nên khởi động lại máy chủ giữa giờ sẽ khiến
        các phiên đang mở phải nhận diện lại từ đầu.
  \item \textbf{Chưa chạy thử ở quy mô thật.} Hệ thống mới chạy trên máy phát
        triển và dữ liệu mẫu, chưa qua thử nghiệm với nhiều lớp học đồng thời
        trên hạ tầng của nhà trường.
\end{enumerate}

\section{Kết luận}

Đồ án đã xây dựng hoàn chỉnh FAMS, một hệ thống quản lý sinh viên tích hợp
điểm danh bằng nhận diện khuôn mặt, chạy được từ đầu tới cuối trên máy chủ
không có GPU. Sản phẩm gồm một dịch vụ với \FsoEndpoint{} endpoint trên
\FsoBang{} bảng dữ liệu, một ứng dụng Flutter \FsoManHinh{} màn hình dùng
chung cho Web, Android và iOS, phân quyền cho \FsoVaiTro{} vai trò theo đúng
cơ cấu tổ chức của trường, cùng \FsoTest{} ca kiểm thử tự động chạy trên
PostgreSQL thật.

Cả năm mục tiêu đặt ra ở Mục~\ref{sec:muctieu} đều đã đạt, với bằng chứng
đối chiếu ở Bảng~\ref{tab:doichieumuctieu}. Phần chưa trọn vẹn nằm ở khâu
kiểm định: hệ thống chưa được đo độ chính xác nhận diện và khả năng chống giả
mạo trên dữ liệu thật, và chưa chạy thử ở quy mô thật của nhà trường.

\section{Đóng góp của đồ án}

\begin{itemize}[leftmargin=1.6em, itemsep=3pt]
  \item \textbf{Một quy trình điểm danh nhiều lớp phòng vệ.} Thay vì coi
        nhận diện khuôn mặt là lời giải duy nhất, hệ thống xếp chồng sáu lớp
        kiểm tra: chất lượng ảnh, chống giả mạo, xác nhận đa khung hình, thu
        hẹp không gian so khớp, đối chiếu vị trí theo phòng học và mã QR xoay
        vòng.
  \item \textbf{Mô hình phân quyền hai chiều.} Quyền theo chức năng kết hợp
        với phạm vi dữ liệu theo Khoa, cài đặt tập trung dưới dạng bộ lọc
        truy vấn, cho phép thêm vai trò mới mà không phải sửa từng endpoint;
        mọi thao tác quản lý đều được ghi vào nhật ký hệ thống.
  \item \textbf{Thiết kế tôn trọng dữ liệu sinh trắc.} Đồng ý sinh trắc học
        là bắt buộc, có phiên bản nội dung và được kiểm tra ở mọi điểm xử lý;
        chỉ lưu vector đặc trưng chứ không lưu ảnh lâu dài, ảnh bằng chứng
        phát qua đường dẫn được ký điện tử và tự hết hạn, người dùng rút đồng
        ý là dữ liệu bị xoá ngay.
  \item \textbf{Bộ tài liệu tự đồng bộ với mã nguồn.} Mọi con số về quy mô
        trong báo cáo này được sinh tự động từ kho mã nguồn bằng công cụ mô
        tả ở mục C của phần Phụ lục, nên báo cáo không thể mô tả sai hệ
        thống mà không bị phát hiện.
\end{itemize}

\section{Hướng phát triển}

\begin{enumerate}[leftmargin=1.8em, itemsep=3pt]
  \item \textbf{Đo độ chính xác và khả năng chống giả mạo.} Thu thập ảnh có
        sự đồng ý của nhiều sinh viên ở nhiều điều kiện ánh sáng và góc chụp,
        tính tỉ lệ chấp nhận nhầm và từ chối nhầm \cite{Jain2004Biometric,
        ISO19795}, chọn ngưỡng nhận diện theo đường cong ROC
        \cite{Fawcett2006ROC} thay vì đặt theo kinh nghiệm; thử chống giả mạo
        bằng ảnh in, màn hình và video, ghi cả số lần chặn nhầm người thật.
  \item \textbf{Chỉ nạp mô hình cần dùng.} Giới hạn bộ \code{buffalo\_l} ở
        phát hiện và trích đặc trưng; phép đo ở Mục~\ref{sec:hieunang} cho
        thấy thời gian giảm \HnGiamMotMat\% với khung một khuôn mặt mà vector
        không đổi.
  \item \textbf{Ngưỡng thích nghi theo từng người.} Tính ngưỡng riêng dựa
        trên độ phân tán của các vector đã đăng ký của mỗi sinh viên, thay
        cho một ngưỡng chung.
  \item \textbf{Chống giả mạo mạnh hơn.} Bổ sung tín hiệu chiều sâu nếu
        thiết bị có cảm biến phù hợp, hoặc kết hợp thêm kiểm tra chủ động
        trong luồng sinh viên tự điểm danh.
  \item \textbf{Đưa bộ đếm xác nhận ra kho lưu trữ ngoài tiến trình.} Dùng
        một kho khoá và giá trị như Redis để phiên đang mở không bị mất
        trạng thái khi máy chủ khởi động lại, đồng thời chạy được nhiều tiến
        trình máy chủ song song.
  \item \textbf{Buộc đổi mật khẩu sau khi được đặt lại.} Hiện cán bộ đặt lại
        xong vẫn biết mật khẩu của người đó. Nên yêu cầu người dùng đổi sang
        mật khẩu riêng ngay ở lần đăng nhập kế tiếp.
  \item \textbf{Kết nối với hệ thống quản lý đào tạo của trường.} Đồng bộ
        danh sách lớp học phần và ghi danh thay vì nhập lại, đây là điều
        kiện bắt buộc nếu muốn triển khai thật.
  \item \textbf{Thông báo đẩy trên điện thoại.} Phụ huynh hiện nhận thông
        báo trong hộp thư của ứng dụng và qua email. Bước tiếp theo là thông
        báo đẩy qua dịch vụ của hệ điều hành, và gửi cảnh báo cho cả sinh viên
        ngay khi tỉ lệ vắng chạm ngưỡng \FtiLeCanhBao\%.
\end{enumerate}
